\documentclass{article}
\usepackage{graphicx} % Required for inserting images
\usepackage{geometry}
\usepackage{algorithm}
\usepackage{algorithmic}
\usepackage{amsmath}
\usepackage{amssymb}
\geometry{a4paper,
left=2cm,
right=2cm,
top=2cm,
bottom=2cm,
}
\usepackage[
backend=biber,
style=ieee,
citestyle=numeric-comp
]{biblatex}
\addbibresource{ProposalReferences.bib}

%-----------------------------------------------------%

\begin{document}
	
	\begin{center}
		\Large Development of a Bayesian Optimization algorithm to explore magnetic field topology for radiation dose reduction in crewed space missions shielding.
	\end{center}
	
	% Things needed/points to hit:
	% WHY is this research NEEDED, what's the biggest picture this influences, who does it help, why should we do it? Be grandiose at first
	% HOW is this reserach going to be different to what's already been done, what previous research has been done if any, where are the gaps and how will this research fill it?
	% How will it be done/ how will the research questions be answered?
	
	The problem of radiation shielding is of paramount importance in crewed space missions: the limitations of damages to the astronauts' healths puts limits on the duration of a crewed mission. At this time, NASA duration for ISS missions is 6 months \cite{Kahn2014}. This cap makes unfeasible at the current stage the possibility of planning crewed mission to explore other solar system planets. 
	
	There are two types of radiation shielding: passive and active. Passive shielding operates by placing some shielding material that absorbs radiation by particle interaction. This is the type of shielding that the most research has been conducted about.
	
	Active radiation shielding relies on a magnetic field to deflect incoming particles via Lorentz force. Some designs have been proposed, like in \cite{Vuolo2016,undergrad2018,Washburn2015}. Also, preliminary studies have been conducted by NASA concluding that using magnetic fields generated by high temperature superconducting coils is possible, while the proposed designs do not offer significant improvements over passive shielding \cite{niac2014nasa}. At the current state, no algorithms for BO in magnetic field optimization have been proposed in the reviewed literature.
	\\
	
	In the current literature, the most widely used method for magnetic field optimization is topology optimization, with applications to finding the best coil distribution for the largest uniform field in \cite{Thabuis2022} and in reverse magnetostatics to find optimal configuration of coils to generate a target field in \cite{KAPTANOGLU2024}. This may prove useful to reverse engineer an optimal solution from Bayesian Optimization. 
	Given the experience of the RMIT Space Physics Group in this challenge, as in \cite{Iles2021,Ng2025,LiamM}, this research proposes using a statistical optimization model to explore exotic field topologies for minimizing the received radiation dose in Monte Carlo simulations.
	Due to the complex and non-linear nature of dependencies between field topology and shielding capabilities, no magnetic field optimization has been applied to this task. However we feel that the huge exploration capabilities, even when searches are expensive, and the ability to operate in a black box scenario of BO may provide an important insight in designing the shield configuration.
	\\
	
	\textbf{Research Questions:}
	\begin{enumerate}
		\item Is there an optimal field topology that minimizes the probability of particle penetration and can it be found via Bayesian Optimization?
		\item What low-dimensional parametrizations of magnetic fields produces a physically meaningful optimum while enabling efficient Bayesian Optimization?
		\item What theoretical guarantees can be established for Gaussian Process–based Bayesian Optimization in this PDE-constrained electromagnetic design problem, where the objective is noisy, non-smooth, and induced by stochastic charged-particle transport?
	\end{enumerate}
	
	\section*{Methods}
	The main mathematical tool is the Gaussian Process–based Bayesian Optimization, which tries to minimize the functional $\mathcal{J}(\theta)$ by observing data from it and without having to know it \cite{frazier2018tutorialbayesianoptimization}. Since the simulation of $\mathcal{J}$ is complex, this avoids non-linearity and non-differentiability issues of classical optimization algorithms.
	The algorithm works as follow:
	\begin{algorithm}
		\caption{Basic BO minimization}
		\begin{algorithmic}
			\STATE Place a Gaussian process prior on $\mathcal{J}$.
			\STATE Observe $\mathcal{J}$ at an $n_0$ number of points and set $n = n_0$.
			\WHILE{$(n \le N)$}
			\STATE Update posterior probability distribution of $\mathcal{J}$, using all available data.
			\STATE Let $\theta_n$ be a maximizer of an acquisition function over $\theta$.
			\STATE Observe $y_n = \mathcal{J}(\theta_n)$.
			\STATE Increment $n$.
			\ENDWHILE
			\STATE Return point with highest evaluated $\mathcal{J}(\theta)$
		\end{algorithmic}
	\end{algorithm}
	
	The first task will be to find a parametrization of the $\mathbf{B}$ field that is both mathematically tractable and physically informative, given that BO works better in low dimension (max 50 params. approx.).
	
	Then the algorithm will be implemented in a python code that will pass the field parameters to a C++ Monte Carlo simulations (the $\mathcal{J}$ functional), that will return the received radiation dose. The RMIT Space Physics group has provided the computational resources.
	
	The RMIT Space Physics group has already developed simulators for this task that run on the particle physics software GEANT4, that can be used for the task.
	
	A theoretical study on convergence guarantees and robustness under noise will be conducted in parallel, with the help of italian supervisors \cite{boconvergence2024}.
	
	\section{Preliminary study on objective function continuity}
	% Define why it is important
	As the literature highlights, the only requirement for Bayesian Optimization is for the objective function to be continuous, otherwise the algorithm may not converge \cite{frazier2018tutorialbayesianoptimization}. We suppose now that a parametrization such that $\mathbf{B}(\theta)$ is continuous w.r.t. $\theta \in \Theta$ exists. This is actually obvious, since we can always construct our parametrization to be continuous, by ignoring the physical meaning. Now the theoretical challenge is to make the radiation dose to depend continuously on the magnetic field. We will tackle this problem by dividing it in two pieces.
	
	\subsection{Trajectory continuity}
	% Trajectory Mapping and Lorentz Force
	The first point is to make sure that the trajectory of a generic charged particle  depends continuously on the magnetic field. The trajectory is defined as the unique solution $\mathbf{x}(t; \mathbf{B}, \xi)$ to the initial value problem governed by the Lorentz force:
	\begin{equation}
		\begin{cases}
			\mathbf{x}(0) = \mathbf{x}_0, \ \dot{\mathbf{x}}(0) = \mathbf{v}_0, \\
			m \frac{d^2\mathbf{x}}{dt^2} = q \left( \frac{d\mathbf{x}}{dt} \times \mathbf{B}(\mathbf{x}) \right),
		\end{cases}
	\end{equation}
	where $\mathbf{B}$ is a magnetic field in the Banach space $C^1(\Theta, \mathbb{R}^3)$ and $\xi = (\mathbf{x}_0, \mathbf{v}_0) \in \Xi$ represents the initial state of a particle in the probability space of incoming radiation.
	
	Let's now rewrite the system as a first order ODE problem. Let $\gamma(t) = (\mathbf{x}(t), \mathbf{v}(t))$ be the trajectory of a particle in the field, then:
	
	\begin{equation}
		\begin{cases}
			\gamma(0) = \xi, \\
			{d \gamma \over dt} = F(\gamma, \mathbf{B}, t),
		\end{cases}
	\end{equation}
	where:
	
	\begin{equation}
		F(\gamma, \mathbf{B}, t) =
		\begin{pmatrix}
			\mathbf{v}(t) \\
			{q \over m} \left( \mathbf{v}(t) \times \mathbf{B}(\mathbf{x}) \right)
		\end{pmatrix}
	\end{equation}
	is the nonlinear bounded functional representing the Lorentz force.
	
	We can exploit the following facts:
	\begin{enumerate}
		\item $\mathbf{B}$ is continuous and Lipschitz by construction;
		\item $\| \mathbf{v} \|$ is physically bounded by the speed of light;
	\end{enumerate}  
	to apply the Picard-Lindelof theorem, which ensures existence and unicity of solution.
	
	\subsection{Revelator continuity}
	% Soft Scoring Kernel Definitions
	To ensure the continuity of the objective function, we need to avoid using a hit/miss revelator, since step function is non continuous. We define a soft scoring kernel $\Phi: \mathbb{R}^3 \to [0, 1]$. A possible choice for ensuring $C^\infty$ regularity is a Normalized Compact Mollifier with support on the ball $\mathcal{B}_\varepsilon$ of radius $\epsilon$:
	
	\begin{equation}
		\Phi(\mathbf{x}) = 
		\begin{cases} 
			{1 \over \lambda_n(\mathcal{B}_\varepsilon)} \exp\left(-\frac{1}{1 - \|\mathbf{x}/\epsilon\|^2}\right) & \text{if } \|\mathbf{x}\| < \epsilon, \\
			0 & \text{otherwise};
		\end{cases}
	\end{equation}
	where $\lambda_n(\mathcal{B}_\varepsilon) = \int_{\mathcal{B}_\varepsilon}  \exp\left(-\frac{1}{1 - \|\mathbf{x}/\epsilon\|^2}\right) d\lambda_n$, with $\lambda_n$ as the $n$-dimensional Lebesgue measure, is the volume of the ball. Note: in this setup, the mollifier is also a probability distribution. This may be useful later.
	
	% Expected Dose Functional
	The individual dose $D(\mathbf{B}, \xi)$ received for a single particle is then defined as the line integral along its trajectory $\gamma(s)$:
	
	\begin{equation}
		D(\mathbf{B}, \xi) = \int_\gamma \Phi(\mathbf{x}(t; \mathbf{B}, \xi)) ds = \int_0^T \Phi(\mathbf{x}(t; \mathbf{B}, \xi)) \cdot \|\mathbf{v}(t; \mathbf{B}, \xi)\| dt.
	\end{equation}
	
	Since Boris integrator should conserve energy \cite{borisalgorithm}, we can eliminate time dependence from the speed norm, as it should be almost constant, thus obtaining:
	
	\begin{equation}
		D(\mathbf{B}, \xi) = \|\mathbf{v}(\xi)\| \int_0^T \Phi(\mathbf{x}(t; \mathbf{B}, \xi)) \ dt.
	\end{equation}
	
	This allows us to reduce calls to \verb|Eigen::Vector3d.norm()|. 
	
	To conclude this discussion, the global objective function to be minimized is the expected received dose $\mathcal{J}(\mathbf{\Theta})$, defined by the expectation over the radiation spectrum $P(\xi)$:
	
	\begin{equation}
		\mathcal{J}(\mathbf{\Theta}) = \mathbb{E}_{\xi \sim P} [D(\mathbf{B}, \xi)] = \int_{\Xi} D(\mathbf{B}, \xi) P(\xi) \ d\xi = \int_{\Xi} \|\mathbf{v}(\xi)\|  P(\xi) \int_0^T \Phi(\mathbf{x}(t; \mathbf{B}, \xi)) \ dt \ d\xi.
	\end{equation}

	We take the expectation of the dose to try to average out the intrinsic variance of the Monte Carlo simulation. At this point, our observations will hopefully come from a continuous function (maybe even smooth) with some small random noise. 
	
	The study of robustness of BO under noise will come in the next chapters, but we already have some guarantees in literature, like in \cite{boconvergence2024}, where Gardner et al. provide bounds for local optimization in high dimension under noisy conditions.
	
	
	\section{Remaining questions}
	There are still a few more question to address for this application of BO:
	\begin{enumerate}
		\item What are the guarantees of convergence under noisy conditions (crucial if the Monte Carlo simulation has high variance)?
		\item How do this problem's convergence rate and regret bound scale with dimension?
		\item What is the most useful parametrization, both from a mathematical and a physical perspective?
	\end{enumerate}
	
	
	Additional references include: \cite{rasmussen_2005}, \cite{garnett_bayesoptbook_2023}.
	
	\clearpage
	\printbibliography
\end{document}
